\section{Developing the Schedule}
\subsection{Activity List, Resources and Durations}
\textbf{Method.} The primary owner of each activity estimated it, drawing on what the work actually involved, so these are retrospective estimates from commit history and team recall, not forecasts made at the time. Hours are shown by member (Su, Ul, Sh). Guide reviews (WBS 7.1) are level of effort and are not in the network.

\begin{longtable}{>{\raggedright\arraybackslash}p{0.65cm}>{\raggedright\arraybackslash}p{6.52cm}>{\raggedright\arraybackslash}p{2.06cm}>{\raggedright\arraybackslash}p{3.59cm}>{\centering\arraybackslash}p{1.41cm}>{\raggedright\arraybackslash}p{2.07cm}}
\hline\rowcolor{navy}\tablelist{Table 4.1 - Activity list with resource and duration estimates}\thead{\#} & \thead{Activity} & \thead{Predecessors} & \thead{Resources (hours)} & \thead{Duration (days)} & \thead{Estimated by} \\ \hline\endfirsthead
\hline\rowcolor{navy}\thead{\#} & \thead{Activity} & \thead{Predecessors} & \thead{Resources (hours)} & \thead{Duration (days)} & \thead{Estimated by} \\ \hline\endhead
A & Literature survey & - & Su 16, Ul 16, Sh 17 & 10 & Sushanth \\ \hline
B & Requirements and architecture & A & Su 20, Ul 8, Sh 9 & 15 & Sushanth \\ \hline
C & Platform scaffold and core & B & Sh 21, Su 8 & 10 & Shreyas \\ \hline
D & Backend API and database & C & Su 33, Sh 4 & 20 & Sushanth \\ \hline
E & Dashboard v1 and redesign & C & Sh 38, Su 4, Ul 4 & 30 & Shreyas \\ \hline
F & Cap dataset, detector and classifier iterations & B & Ul 44, Sh 8 & 25 & Ullas \\ \hline
G & Fill dataset and fill-level models & F & Ul 20, Sh 8 & 10 & Ullas \\ \hline
H & First integrated demo & D, E, G & Ul 8, Su 4, Sh 4 & 5 & Ullas \\ \hline
I & Thesis, IEEE paper, presentation & H & Su 20, Ul 20, Sh 8 & 10 & Sushanth \\ \hline
J & Label: PatchCore baseline & I & Ul 28, Sh 4 & 25 & Ullas \\ \hline
K & Label: annotation, YOLO iteration, template matching & J & Sh 25, Ul 20 & 15 & Shreyas \\ \hline
L & Label: EfficientAD, grayscale, calibration & K & Ul 36, Su 4 & 15 & Ullas \\ \hline
M & Label dashboard & L & Sh 25 & 15 & Shreyas \\ \hline
N & Cap teacher-student distillation & L & Su 25, Ul 4 & 10 & Sushanth \\ \hline
O & Station 1 production pipeline & N & Su 37, Sh 4 & 10 & Sushanth \\ \hline
P & Docker and TensorRT deployment & O & Su 29 & 8 & Sushanth \\ \hline
Q & Audit, Replay, Active Learning workstations & O & Su 29 & 10 & Sushanth \\ \hline
R & Station 2 integration, env-driven paths & P, Q & Su 21 & 5 & Sushanth \\ \hline
S & GPU verification and handover & M, R & Ul 8, Su 4, Sh 4 & 3 & Ullas \\ \hline
\end{longtable}\nopagebreak
\vspace{-1.2em}\nopagebreak\tcap{Table 4.1 - Activity list with resource and duration estimates}


\subsection{Network Diagram and Schedule Calculations}
The network diagram is in Appendix B. Days are working days counted from 0 at the start of 2 February 2026.

\begin{longtable}{>{\raggedright\arraybackslash}p{0.61cm}>{\raggedright\arraybackslash}p{6.93cm}>{\centering\arraybackslash}p{0.91cm}>{\centering\arraybackslash}p{0.91cm}>{\centering\arraybackslash}p{0.91cm}>{\centering\arraybackslash}p{0.91cm}>{\centering\arraybackslash}p{0.91cm}>{\centering\arraybackslash}p{0.91cm}>{\centering\arraybackslash}p{0.91cm}>{\centering\arraybackslash}p{1.59cm}}
\hline\rowcolor{navy}\tablelist{Table 4.2 - Forward and backward pass, slack and critical path}\thead{\#} & \thead{Activity} & \thead{Dur.} & \thead{ES} & \thead{EF} & \thead{LS} & \thead{LF} & \thead{Total slack} & \thead{Free slack} & \thead{Critical?} \\ \hline\endfirsthead
\hline\rowcolor{navy}\thead{\#} & \thead{Activity} & \thead{Dur.} & \thead{ES} & \thead{EF} & \thead{LS} & \thead{LF} & \thead{Total slack} & \thead{Free slack} & \thead{Critical?} \\ \hline\endhead
A & Literature survey & 10 & 0 & 10 & 0 & 10 & 0 & 0 & \textbf{Yes} \\ \hline
B & Requirements and architecture & 15 & 10 & 25 & 10 & 25 & 0 & 0 & \textbf{Yes} \\ \hline
C & Platform scaffold and core & 10 & 25 & 35 & 25 & 35 & 0 & 0 & \textbf{Yes} \\ \hline
D & Backend API and database & 20 & 35 & 55 & 45 & 65 & 10 & 10 & No \\ \hline
E & Dashboard v1 and redesign & 30 & 35 & 65 & 35 & 65 & 0 & 0 & \textbf{Yes} \\ \hline
F & Cap dataset, detector and classifier iterations & 25 & 25 & 50 & 30 & 55 & 5 & 0 & No \\ \hline
G & Fill dataset and fill-level models & 10 & 50 & 60 & 55 & 65 & 5 & 5 & No \\ \hline
H & First integrated demo & 5 & 65 & 70 & 65 & 70 & 0 & 0 & \textbf{Yes} \\ \hline
I & Thesis, IEEE paper, presentation & 10 & 70 & 80 & 70 & 80 & 0 & 0 & \textbf{Yes} \\ \hline
J & Label: PatchCore baseline & 25 & 80 & 105 & 80 & 105 & 0 & 0 & \textbf{Yes} \\ \hline
K & Label: annotation, YOLO iteration, template matching & 15 & 105 & 120 & 105 & 120 & 0 & 0 & \textbf{Yes} \\ \hline
L & Label: EfficientAD, grayscale, calibration & 15 & 120 & 135 & 120 & 135 & 0 & 0 & \textbf{Yes} \\ \hline
M & Label dashboard & 15 & 135 & 150 & 155 & 170 & 20 & 20 & No \\ \hline
N & Cap teacher-student distillation & 10 & 135 & 145 & 135 & 145 & 0 & 0 & \textbf{Yes} \\ \hline
O & Station 1 production pipeline & 10 & 145 & 155 & 145 & 155 & 0 & 0 & \textbf{Yes} \\ \hline
P & Docker and TensorRT deployment & 8 & 155 & 163 & 157 & 165 & 2 & 2 & No \\ \hline
Q & Audit, Replay, Active Learning workstations & 10 & 155 & 165 & 155 & 165 & 0 & 0 & \textbf{Yes} \\ \hline
R & Station 2 integration, env-driven paths & 5 & 165 & 170 & 165 & 170 & 0 & 0 & \textbf{Yes} \\ \hline
S & GPU verification and handover & 3 & 170 & 173 & 170 & 173 & 0 & 0 & \textbf{Yes} \\ \hline
\end{longtable}\nopagebreak
\vspace{-1.2em}\nopagebreak\tcap{Table 4.2 - Forward and backward pass, slack and critical path}


\textbf{Critical path:} \vCriticalPath{}, totalling \vProjectDur{} working days.

\textbf{Worked calculation for activity D (Backend API and database, 20 days).}
\begin{itemize}
\item Forward pass, Rule 1: D's only predecessor is C, with EF 35, so ES = 35 and EF = 35 + 20 = 55.
\item Backward pass, Rule 2: D's only successor is H, with LS 65, so LF = 65 and LS = 65 - 20 = 45.
\item Total slack = LS - ES = 45 - 35 = 10 days.
\item Free slack = earliest successor ES - EF = 65 - 55 = 10 days.
\end{itemize}

\textbf{(a)} Activity \vSlackMostId{}, \vSlackMostName{}, has the most total slack at \vSlackMost{} days. We would use it to lend Shreyas to critical-path work when behind, and indeed the dashboard consumed 15 of those days without moving the finish.

\textbf{(b)} The largest critical activity is \vLongestId{}, \vLongestName{}, at \vLongestDur{} days. A five-day slip would move the finish from day 173 to day 178, 7 October 2026, making the total slack on the critical path -5 against the required date. Paths joining E at H gain slack against the later finish: D rises from 10 to 15 days, and F and G from 5 to 10. Paths after H keep theirs, because their dates all move together: M stays at 20 days and P at 2.

\subsection{Milestones}
\begin{longtable}{>{\raggedright\arraybackslash}p{8.41cm}>{\raggedright\arraybackslash}p{2.43cm}>{\raggedright\arraybackslash}p{3.1cm}>{\raggedright\arraybackslash}p{2.76cm}}
\hline\rowcolor{navy}\thead{Milestone} & \thead{Planned date} & \thead{Actual date} & \thead{Status} \\ \hline\endfirsthead
\hline\rowcolor{navy}\thead{Milestone} & \thead{Planned date} & \thead{Actual date} & \thead{Status} \\ \hline\endhead
Literature survey complete & 13 Feb & 16 Feb & Late by 1 day \\ \hline
Requirements and architecture complete & 6 Mar & 6 Mar & Met \\ \hline
Platform scaffold pushed & 20 Mar & 21 Mar & Met \\ \hline
Cap models accepted & 10 Apr & 3 May & Late by 15 days \\ \hline
First integrated demo & 8 May & 15 May & Late by 5 days \\ \hline
Thesis complete & 22 May & 22 May & Met \\ \hline
Label models calibrated & 7 Aug & 7 Aug & Met \\ \hline
Student cap model accepted & 21 Aug & 7 Sep & Late by 11 days \\ \hline
Station 1 production pipeline & 4 Sep & 14 Sep & Late by 6 days \\ \hline
Station 2 integrated & 25 Sep & 26 Oct & Late by 21 days \\ \hline
GPU verification and handover & 30 Sep & 26 Oct & Late by 18 days \\ \hline
\end{longtable}\nopagebreak
\vspace{-1.2em}\nopagebreak\tcap{Table 4.3 - Milestones, all in 2026 (a milestone has zero duration)}


\subsection{Baseline Schedule}
\begin{center}\resizebox{\textwidth}{!}{\begin{tikzpicture}[font=\scriptsize]
\draw[gridline] (0.00,0.25) -- (0.00,-8.74);
\node[anchor=south] at (0.56,0.2) {Feb};
\draw[gridline] (1.15,0.25) -- (1.15,-8.74);
\node[anchor=south] at (1.77,0.2) {Mar};
\draw[gridline] (2.43,0.25) -- (2.43,-8.74);
\node[anchor=south] at (3.03,0.2) {Apr};
\draw[gridline] (3.66,0.25) -- (3.66,-8.74);
\node[anchor=south] at (4.28,0.2) {May};
\draw[gridline] (4.94,0.25) -- (4.94,-8.74);
\node[anchor=south] at (5.54,0.2) {Jun};
\draw[gridline] (6.18,0.25) -- (6.18,-8.74);
\node[anchor=south] at (6.79,0.2) {Jul};
\draw[gridline] (7.45,0.25) -- (7.45,-8.74);
\node[anchor=south] at (8.07,0.2) {Aug};
\draw[gridline] (8.73,0.25) -- (8.73,-8.74);
\node[anchor=south] at (9.33,0.2) {Sep};
\draw[gridline] (9.96,0.25) -- (9.96,-8.74);
\node[anchor=south] at (10.58,0.2) {Oct};
\draw[gridline] (11.20,0.25) -- (11.20,-8.74);
\node[anchor=east] at (-0.1,-0.23) {A\ \ Literature survey};
\fill[critbar] (0.04,-0.36) rectangle (0.54,-0.10);
\node[anchor=east] at (-0.1,-0.69) {B\ \ Requirements and architecture};
\fill[critbar] (0.62,-0.82) rectangle (1.40,-0.56);
\node[anchor=east] at (-0.1,-1.15) {C\ \ Platform scaffold and core};
\fill[critbar] (1.48,-1.28) rectangle (1.98,-1.02);
\node[anchor=east] at (-0.1,-1.61) {D\ \ Backend API and database};
\fill[plainbar] (2.06,-1.74) rectangle (3.13,-1.48);
\draw[slackline] (3.13,-1.61) -- (3.71,-1.61);
\draw[slackline,solid] (3.71,-1.68) -- (3.71,-1.54);
\node[anchor=east] at (-0.1,-2.07) {E\ \ Dashboard v1 and redesign};
\fill[critbar] (2.06,-2.20) rectangle (3.71,-1.94);
\node[anchor=east] at (-0.1,-2.53) {F\ \ Cap dataset, detector and classifier iterations};
\fill[plainbar] (1.48,-2.66) rectangle (2.84,-2.40);
\draw[slackline] (2.84,-2.53) -- (3.13,-2.53);
\draw[slackline,solid] (3.13,-2.60) -- (3.13,-2.46);
\node[anchor=east] at (-0.1,-2.99) {G\ \ Fill dataset and fill-level models};
\fill[plainbar] (2.92,-3.12) rectangle (3.42,-2.86);
\draw[slackline] (3.42,-2.99) -- (3.71,-2.99);
\draw[slackline,solid] (3.71,-3.06) -- (3.71,-2.92);
\node[anchor=east] at (-0.1,-3.45) {H\ \ First integrated demo};
\fill[critbar] (3.79,-3.58) rectangle (3.99,-3.32);
\node[anchor=east] at (-0.1,-3.91) {I\ \ Thesis, IEEE paper, presentation};
\fill[critbar] (4.08,-4.04) rectangle (4.57,-3.78);
\node[anchor=east] at (-0.1,-4.37) {J\ \ Label: PatchCore baseline};
\fill[critbar] (4.65,-4.50) rectangle (6.01,-4.24);
\node[anchor=east] at (-0.1,-4.83) {K\ \ Label: annotation, YOLO iteration, template matching};
\fill[critbar] (6.09,-4.96) rectangle (6.88,-4.70);
\node[anchor=east] at (-0.1,-5.29) {L\ \ Label: EfficientAD, grayscale, calibration};
\fill[critbar] (6.96,-5.42) rectangle (7.74,-5.16);
\node[anchor=east] at (-0.1,-5.75) {M\ \ Label dashboard};
\fill[plainbar] (7.82,-5.88) rectangle (8.61,-5.62);
\draw[slackline] (8.61,-5.75) -- (9.76,-5.75);
\draw[slackline,solid] (9.76,-5.82) -- (9.76,-5.68);
\node[anchor=east] at (-0.1,-6.21) {N\ \ Cap teacher-student distillation};
\fill[critbar] (7.82,-6.34) rectangle (8.32,-6.08);
\node[anchor=east] at (-0.1,-6.67) {O\ \ Station 1 production pipeline};
\fill[critbar] (8.40,-6.80) rectangle (8.89,-6.54);
\node[anchor=east] at (-0.1,-7.13) {P\ \ Docker and TensorRT deployment};
\fill[plainbar] (8.98,-7.26) rectangle (9.39,-7.00);
\draw[slackline] (9.39,-7.13) -- (9.47,-7.13);
\draw[slackline,solid] (9.47,-7.20) -- (9.47,-7.06);
\node[anchor=east] at (-0.1,-7.59) {Q\ \ Audit, Replay, Active Learning workstations};
\fill[critbar] (8.98,-7.72) rectangle (9.47,-7.46);
\node[anchor=east] at (-0.1,-8.05) {R\ \ Station 2 integration, env-driven paths};
\fill[critbar] (9.55,-8.18) rectangle (9.76,-7.92);
\node[anchor=east] at (-0.1,-8.51) {S\ \ GPU verification and handover};
\fill[critbar] (9.84,-8.64) rectangle (9.96,-8.38);
\draw[marker] (9.96,0) -- (9.96,-8.74);
\node[anchor=north,align=center] at (9.96,-8.74) {Planned finish\\30 Sep};
\fill[critbar] (-5.40,-9.69) rectangle (-4.90,-9.51);
\node[anchor=west] at (-4.85,-9.60) {Critical};
\fill[plainbar] (-3.40,-9.69) rectangle (-2.90,-9.51);
\node[anchor=west] at (-2.85,-9.60) {Non-critical};
\draw[slackline] (-0.60,-9.60) -- (-0.10,-9.60);
\node[anchor=west] at (-0.05,-9.60) {Total slack};
\end{tikzpicture}
}\end{center}
\tcap{Figure 4.1 - Baseline schedule, critical path in orange, total slack dashed}

Planned finish and required completion date are both 30 September 2026, so total slack on the critical path is zero.

\subsection{Project Control Process}
The reporting period is one calendar month, matched to how the work actually ran alongside coursework; guide reviews took place about weekly. Data for each period are actual start and finish (from commits), hours, costs expended and committed (always ₹0 committed) and percent complete. The team and the guide attend reviews. The threshold we set when formalising this process is a slip of more than 5 working days on a critical activity, or a CPI below 0.90.

We did not run this process formally at the time: the baseline was formalised retrospectively for this report, and the monthly data in Appendix C are reconstructed from Git. Real control happened through discussion among the team and the guide's reviews.

\subsection{Effects of Actual Schedule Performance}
Activities not listed finished on their planned date.

\begin{longtable}{>{\raggedright\arraybackslash}p{3.57cm}>{\raggedright\arraybackslash}p{1.3cm}>{\raggedright\arraybackslash}p{1.3cm}>{\centering\arraybackslash}p{1.46cm}>{\raggedright\arraybackslash}p{4.98cm}>{\raggedright\arraybackslash}p{3.69cm}}
\hline\rowcolor{navy}\thead{Activity} & \thead{Planned finish} & \thead{Actual finish} & \thead{Variance (days)} & \thead{Cause} & \thead{Effect on total slack} \\ \hline\endfirsthead
\hline\rowcolor{navy}\thead{Activity} & \thead{Planned finish} & \thead{Actual finish} & \thead{Variance (days)} & \thead{Cause} & \thead{Effect on total slack} \\ \hline\endhead
A Literature survey & 13 Feb & 16 Feb & +1 & Survey ran one day past its two weeks & -1 on the critical path; recovered within B \\ \hline
E Dashboard v1 and redesign & 1 May & 4 May & +1 & Redesign landed on 4 May & -1 on the critical path \\ \hline
F Cap dataset, detector and classifier iterations & 10 Apr & 3 May & +15 & Bottle cap models needed repeated iteration: small caps, class mapping, white-only colour bias & Used all 5 days of slack; path at -10 \\ \hline
G Fill dataset and fill-level models & 24 Apr & 15 May & +15 & Inherited F's delay; water datasets had to be merged & Path stayed at -10 \\ \hline
H First integrated demo & 8 May & 15 May & +5 & Waited on G & -5 on the critical path; recovered by overlapping I \\ \hline
J Label: PatchCore baseline & 26 Jun & 11 Jul & +10 & Paper revision and a pause after the thesis delayed the start to mid-June & -10 on the critical path \\ \hline
K Label: annotation, YOLO iteration, template matching & 17 Jul & 4 Aug & +12 & YOLO ROI detector tried, then replaced by template matching & -12 on the critical path \\ \hline
M Label dashboard & 28 Aug & 18 Sep & +15 & Label dashboard finished after the models & Used 15 of its 20 days of slack \\ \hline
N Cap teacher-student distillation & 21 Aug & 7 Sep & +11 & CSITSS rejection (22 Aug) took priority, delaying the start & -11 on the critical path \\ \hline
O Station 1 production pipeline & 4 Sep & 14 Sep & +6 & Inherited N's delay & -6 on the critical path \\ \hline
P Docker and TensorRT deployment & 16 Sep & 22 Oct & +26 & TensorRT static batch mismatch; dynamic batching and versioning reworked & Used its 2 days of slack; path at -24 \\ \hline
Q Audit, Replay, Active Learning workstations & 18 Sep & 22 Oct & +24 & Started after the rework on P began; delivered with P on 22 Oct & -24 on the critical path \\ \hline
R Station 2 integration, env-driven paths & 25 Sep & 26 Oct & +21 & Waited on P and Q; Station 2 merge into one UI & -21 on the critical path \\ \hline
S GPU verification and handover & 30 Sep & 26 Oct & +18 & Waited on R & Project finished 18 days late \\ \hline
\end{longtable}\nopagebreak
\vspace{-1.2em}\nopagebreak\tcap{Table 4.4 - Schedule variance analysis}


The early phase recovered by overlapping the thesis with the demo work, so I finished on 22 May as planned, and L also recovered by overlapping K. The later slips did not recover because the distillation and deployment activities depend strictly on each other.

\subsection{Incorporating Changes and Updating the Schedule}
No change forms were used at the time; this log records the decisions visible in the repository, with impact estimated retrospectively.

\simpletable{L{1.11cm}L{6.04cm}L{4.63cm}L{2.52cm}L{2.2cm}}{\thead{ID} & \thead{Description} & \thead{Impact (scope / days / ₹)} & \thead{Decision} & \thead{Approved by}}{
CR1 & Narrow from multi-product to bottle and label lines & Scope reduced; avoided work outside the baseline & Approved & Team, guide \\ \hline
CR2 & Replace PatchCore with per-region EfficientAD & No scope change; part of 10-day J slip; +\vOverJ{} & Approved & Team \\ \hline
CR3 & Replace YOLO localisation with template matching & No scope change; 12-day K slip; +\vOverK{} & Approved & Team \\ \hline
CR4 & Keep label inspection on its own dashboard & Removes merge work; no date impact & Approved & Team \\ \hline
CR5 & Distil YOLOv11x teacher into YOLOv11m student & Inside N; 11-day slip; +\vOverN{} & Approved & Team \\ \hline
CR6 & Rework TensorRT to dynamic batching and versioned engines & Quality up; 26-day P slip; +\vOverP{} & Approved & PM \\ \hline
}{Table 4.5 - Change request log}

\textbf{Updated schedule.} The final finish date is 26 October 2026, 18 working days after the baseline. The revised critical path in the final phase ran through P and Q (both finishing 22 October), then R, ending with S on 26 October. Measured against the 30 September requirement, total slack on that path ended at -18 days.
